\chapter{Useful Matrix Identities}

\vfill
\localtableofcontents
\vfill
\newpage

\section{Gaussian Conditioning Formula}

\begin{proposition}[Gaussian conditioning \cite{bib:multivariate-normal}]\label{th:Gaussian_conditioning}
	For $n,p\in\mathbb{N}$, let $U\in\mathbb{R}^n$ and $V\in\mathbb{R}^p$ be jointly Gaussian random variables, of respective means $\mu_U,\mu_V$ and covariance matrices $\mathbf{A}\succ0$ and $\mathbf{D}\succ0$. There exists $\mathbf{B}\in\mathcal{M}_{n,p}(\mathbb{R})$ such that:
	\begin{equation}
		\begin{bmatrix}U\\V\end{bmatrix}\sim\mathcal{N}\left(\begin{bmatrix}\mu_U\\\mu_V\end{bmatrix},\begin{bmatrix}\mathbf{A}&\mathbf{B}\\\mathbf{B}^\top&\mathbf{D}\end{bmatrix}\right)
	\end{equation}

	We assume $\mathbf{B}$ to be such that this joint covariance is SPD in the sense of \Cref{th:matrix_SPSD}. The conditional distribution of $U$ given $V=y$ is then:
	\begin{equation}
		U|V=y\sim\mathcal{N}\left(\mu_U+\mathbf{B}\mathbf{D}^{-1}(y-\mu_V),\mathbf{P}_{U|V}\right)
	\end{equation}
	with:
	\begin{equation}
		\mathbf{P}_{U|V}\triangleq \mathbf{A}-\mathbf{B}\mathbf{D}^{-1}\mathbf{B}^\top\succ0
	\end{equation}
\end{proposition}

\begin{proof}
	The existence of $\mathbf{B}$ is direct by definition of jointly Gaussian random variables. As $\mathbf{D}\succ0$ is invertible, $\mathbf{S}\triangleq \mathbf{A}-\mathbf{B}\mathbf{D}^{-1}\mathbf{B}^\top$, the Schur complement with respect to $\mathbf{D}$, is well-defined.

	Let $u\neq0$:
	\begin{equation}
		\begin{bmatrix}u\\-\mathbf{D}^{-1}\mathbf{B}^\top u\end{bmatrix}^\top\begin{bmatrix}\mathbf{A}&\mathbf{B}\\\mathbf{B}^\top&\mathbf{D}\end{bmatrix}\begin{bmatrix}u\\-\mathbf{D}^{-1}\mathbf{B}^\top u\end{bmatrix}=u^\top\left(\mathbf{A}-\mathbf{B}\mathbf{D}^{-1}\mathbf{B}^\top\right)u>0
	\end{equation}
	by SPDness of the joint covariance, the left-hand vector being non-zero. $\mathbf{S}$ is thus also invertible, and we can apply the inferior right Schur inversion formula (see \Cref{th:Schur}) to obtain:
	\begin{equation}
		\begin{aligned}
			\begin{bmatrix}\bar{\mathbf{A}}&\bar{\mathbf{B}}\\\bar{\mathbf{C}}&\bar{\mathbf{D}}\end{bmatrix}&\triangleq\begin{bmatrix}\mathbf{A}&\mathbf{B}\\\mathbf{B}^\top&\mathbf{D}\end{bmatrix}^{-1}\\
			&=\begin{bmatrix}\mathbf{S}^{-1}&-\mathbf{S}^{-1}\mathbf{B}\mathbf{D}^{-1}\\-\mathbf{D}^{-1}\mathbf{B}^\top \mathbf{S}^{-1}&\mathbf{D}^{-1}+\mathbf{D}^{-1}\mathbf{B}^\top \mathbf{S}^{-1}\mathbf{B}\mathbf{D}^{-1}\end{bmatrix}
			\end{aligned}
	\end{equation}

	For some realisation $u$ of $U$ and $y$ of $V$, we define $\tilde{u}\triangleq u-\mu_U$ and $\tilde{y}\triangleq y-\mu_V$, and we compute:
	\begin{equation}
		\begin{aligned}
		\alpha&\triangleq\begin{bmatrix}\tilde{u}\\\tilde{y}\end{bmatrix}^\top\begin{bmatrix}\bar{\mathbf{A}}&\bar{\mathbf{B}}\\\bar{\mathbf{C}}&\bar{\mathbf{D}}\end{bmatrix}\begin{bmatrix}\tilde{u}\\\tilde{y}\end{bmatrix}\\
		&=\tilde{u}^\top\bar{\mathbf{A}}\tilde{u}+\tilde{u}^\top\bar{\mathbf{B}}\tilde{y}+\tilde{y}^\top\bar{\mathbf{C}}\tilde{u}+\tilde{y}^\top\bar{\mathbf{D}}\tilde{y}\\
		&=\tilde{u}^\top \mathbf{S}^{-1}\tilde{u}+\tilde{u}^\top \mathbf{S}^{-1}\left(-\mathbf{B}\mathbf{D}^{-1}\tilde{y}\right)+\left(-\mathbf{B}\mathbf{D}^{-1}\tilde{y}\right)^\top \mathbf{S}^{-1}\tilde{u}+\\
		&\qquad\left(-\mathbf{B}\mathbf{D}^{-1}\tilde{y}\right)^\top \mathbf{S}^{-1}\left(-\mathbf{B}\mathbf{D}^{-1}\tilde{y}\right)+\tilde{y}^\top \mathbf{D}^{-1}\tilde{y}\\
		&=(\tilde{u}-\mathbf{B}\mathbf{D}^{-1}\tilde{y})^\top \mathbf{S}^{-1}(\tilde{u}-\mathbf{B}\mathbf{D}^{-1}\tilde{y})+\tilde{y}^\top \mathbf{D}^{-1}\tilde{y}
		\end{aligned}
	\end{equation}

	The conditional density is thus:
	\begin{equation}
		\begin{aligned}
		p\left(\tilde{u}|\tilde{y}\right)&\propto\frac{\displaystyle p\left(\tilde{u},\tilde{y}\right)}{\displaystyle p\left(\tilde{y}\right)}\\
		&\propto\frac{\displaystyle\exp\left(-\frac{1}{2}\alpha\right)}{\displaystyle\exp\left(-\frac{1}{2}\tilde{y}^\top \mathbf{D}^{-1}\tilde{y}\right)}\\
		&\propto\exp\left(-\frac{1}{2}(\tilde{u}-\mathbf{B}\mathbf{D}^{-1}\tilde{y})^\top \mathbf{S}^{-1}(\tilde{u}-\mathbf{B}\mathbf{D}^{-1}\tilde{y})\right)
		\end{aligned}
	\end{equation}

	With this formulation, it comes that $U|V=y$ is Gaussian with mean:
	\begin{equation}
		\mu_{U|V}\triangleq\mu_U+\mathbf{B}\mathbf{D}^{-1}\left(y-\mu_V\right)
	\end{equation}
	and covariance:
	\begin{equation}
		\mathbf{P}_{U|V}\triangleq \mathbf{S}=\mathbf{A}-\mathbf{B}\mathbf{D}^{-1}\mathbf{B}^\top\succ0
	\end{equation}
\end{proof}

\section{Schur Inversion Formula}

\begin{proposition}[Schur inversion formula \cite{bib:Schur-complement}]\label{th:Schur}
	For $n,p\in\mathbb{N}$, let $\mathbf{A}\in\mathcal{M}_{n,n}(\mathbb{R})$, $\mathbf{B}\in\mathcal{M}_{n,p}(\mathbb{R})$, $\mathbf{C}\in\mathcal{M}_{p,n}(\mathbb{R})$, and $\mathbf{D}\in\mathcal{M}_{p,p}(\mathbb{R})$, with $\mathbf{A}$ invertible. We define the matrix $\mathbf{M}$ and its Schur complement with respect to $\mathbf{A}$ as:
	\begin{equation}
		\mathbf{M}\triangleq\begin{bmatrix}\mathbf{A}&\mathbf{B}\\\mathbf{C}&\mathbf{D}\end{bmatrix},\qquad \mathbf{M}/\mathbf{A}\triangleq \mathbf{D}-\mathbf{C}\mathbf{A}^{-1}\mathbf{B}
	\end{equation}

	If $\mathbf{M}/\mathbf{A}$ is invertible, then $\mathbf{M}$ is invertible and:
	\begin{equation}
		\mathbf{M}^{-1}=\begin{bmatrix}\mathbf{A}^{-1}+\mathbf{A}^{-1}\mathbf{B}(\mathbf{M}/\mathbf{A})^{-1}\mathbf{C}\mathbf{A}^{-1}&-\mathbf{A}^{-1}\mathbf{B}(\mathbf{M}/\mathbf{A})^{-1}\\-(\mathbf{M}/\mathbf{A})^{-1}\mathbf{C}\mathbf{A}^{-1}&(\mathbf{M}/\mathbf{A})^{-1}\end{bmatrix}
	\end{equation}
\end{proposition}

\begin{proof}
	Let $\bar{\mathbf{M}}$ be a matrix such that $\mathbf{M}\bar{\mathbf{M}}=\mathbf{I}$, and write its blocks as:
	\begin{equation}
		\begin{bmatrix}\bar{\mathbf{A}}&\bar{\mathbf{B}}\\\bar{\mathbf{C}}&\bar{\mathbf{D}}\end{bmatrix}=\bar{\mathbf{M}}
	\end{equation}

	The identity $\mathbf{M}\bar{\mathbf{M}}=\mathbf{I}$ yields:
	\begin{equation}
		\left\{\begin{aligned}
		\mathbf{A}\bar{\mathbf{A}}+\mathbf{B}\bar{\mathbf{C}}&=\mathbf{I}\\
		\mathbf{A}\bar{\mathbf{B}}+\mathbf{B}\bar{\mathbf{D}}&=0\\
		\mathbf{C}\bar{\mathbf{A}}+\mathbf{D}\bar{\mathbf{C}}&=0\\
		\mathbf{C}\bar{\mathbf{B}}+\mathbf{D}\bar{\mathbf{D}}&=\mathbf{I}\\
		\end{aligned}\right.
	\end{equation}

	In particular, the second equation can be rewritten as:
	\begin{equation}
		\bar{\mathbf{B}}=-\mathbf{A}^{-1}\mathbf{B}\bar{\mathbf{D}}\label{eq:bar_B}
	\end{equation}
	which can be plugged into the fourth equation to give:
	\begin{equation}
		\begin{aligned}
		&\mathbf{C}\bar{\mathbf{B}}+\mathbf{D}\bar{\mathbf{D}}=\mathbf{I}\\
		\iff\quad&-\mathbf{C}\mathbf{A}^{-1}\mathbf{B}\bar{\mathbf{D}}+\mathbf{D}\bar{\mathbf{D}}=\mathbf{I}\\
		\iff\quad&(\mathbf{D}-\mathbf{C}\mathbf{A}^{-1}\mathbf{B})\bar{\mathbf{D}}=\mathbf{I}\\
		\iff\quad&(\mathbf{M}/\mathbf{A})\bar{\mathbf{D}}=\mathbf{I}\\
		\iff\quad&\bar{\mathbf{D}}=(\mathbf{M}/\mathbf{A})^{-1}
		\end{aligned}
	\end{equation}

	\eqref{eq:bar_B} can then be solved:
	\begin{equation}
		\bar{\mathbf{B}}=-\mathbf{A}^{-1}\mathbf{B}(\mathbf{M}/\mathbf{A})^{-1}
	\end{equation}

	A similar operation can be performed on the two remaining equations to obtain:
	\begin{equation}
		\bar{\mathbf{A}}=\mathbf{A}^{-1}(\mathbf{I}-\mathbf{B}\bar{\mathbf{C}})
	\end{equation}
	and:
	\begin{equation}
		\begin{aligned}
		&\mathbf{C}\bar{\mathbf{A}}+\mathbf{D}\bar{\mathbf{C}}=0\\
		\iff\quad&\mathbf{C}\mathbf{A}^{-1}(\mathbf{I}-\mathbf{B}\bar{\mathbf{C}})+\mathbf{D}\bar{\mathbf{C}}=0\\
		\iff\quad&\mathbf{C}\mathbf{A}^{-1}+(\mathbf{D}-\mathbf{C}\mathbf{A}^{-1}\mathbf{B})\bar{\mathbf{C}}=0\\
		\iff\quad&\mathbf{C}\mathbf{A}^{-1}+(\mathbf{M}/\mathbf{A})\bar{\mathbf{C}}=0\\
		\iff\quad&\bar{\mathbf{C}}=-(\mathbf{M}/\mathbf{A})^{-1}\mathbf{C}\mathbf{A}^{-1}
		\end{aligned}
	\end{equation}

	Finally:
	\begin{equation}
		\bar{\mathbf{A}}=\mathbf{A}^{-1}(\mathbf{I}-\mathbf{B}\bar{\mathbf{C}})=\mathbf{A}^{-1}+\mathbf{A}^{-1}\mathbf{B}(\mathbf{M}/\mathbf{A})^{-1}\mathbf{C}\mathbf{A}^{-1}
	\end{equation}

	Conversely, the four blocks of $\bar{\mathbf{M}}$ satisfy:
	\begin{equation}
		\left\{\begin{aligned}
		\mathbf{A}\bar{\mathbf{A}}+\mathbf{B}\bar{\mathbf{C}}&=\mathbf{I}+\mathbf{B}(\mathbf{M}/\mathbf{A})^{-1}\mathbf{C}\mathbf{A}^{-1}-\mathbf{B}(\mathbf{M}/\mathbf{A})^{-1}\mathbf{C}\mathbf{A}^{-1}=\mathbf{I}\\
		\mathbf{A}\bar{\mathbf{B}}+\mathbf{B}\bar{\mathbf{D}}&=-\mathbf{B}(\mathbf{M}/\mathbf{A})^{-1}+\mathbf{B}(\mathbf{M}/\mathbf{A})^{-1}=0\\
		\mathbf{C}\bar{\mathbf{A}}+\mathbf{D}\bar{\mathbf{C}}&=\mathbf{C}\mathbf{A}^{-1}-(\mathbf{D}-\mathbf{C}\mathbf{A}^{-1}\mathbf{B})(\mathbf{M}/\mathbf{A})^{-1}\mathbf{C}\mathbf{A}^{-1}=\mathbf{C}\mathbf{A}^{-1}-\mathbf{C}\mathbf{A}^{-1}=0\\
		\mathbf{C}\bar{\mathbf{B}}+\mathbf{D}\bar{\mathbf{D}}&=-\mathbf{C}\mathbf{A}^{-1}\mathbf{B}(\mathbf{M}/\mathbf{A})^{-1}+\mathbf{D}(\mathbf{M}/\mathbf{A})^{-1}=(\mathbf{D}-\mathbf{C}\mathbf{A}^{-1}\mathbf{B})(\mathbf{M}/\mathbf{A})^{-1}=\mathbf{I}
		\end{aligned}\right.
	\end{equation}
	so that $\mathbf{M}\bar{\mathbf{M}}=\mathbf{I}$. Since $\mathbf{M}$ is a square matrix, it is invertible with inverse $\bar{\mathbf{M}}$.
\end{proof}

\begin{remark}[Inversion with respect to $\mathbf{D}$]
	If $\mathbf{D}$ is invertible, one can instead introduce the Schur complement with respect to $\mathbf{D}$:
	\begin{equation}
		\mathbf{M}/\mathbf{D}\triangleq \mathbf{A}-\mathbf{B}\mathbf{D}^{-1}\mathbf{C}
	\end{equation}

	When $\mathbf{M}/\mathbf{D}$ is also invertible, an analogous block inversion formula holds:
	\begin{equation}
		\mathbf{M}^{-1}=\begin{bmatrix}(\mathbf{M}/\mathbf{D})^{-1}&-(\mathbf{M}/\mathbf{D})^{-1}\mathbf{B}\mathbf{D}^{-1}\\-\mathbf{D}^{-1}\mathbf{C}(\mathbf{M}/\mathbf{D})^{-1}&\mathbf{D}^{-1}+\mathbf{D}^{-1}\mathbf{C}(\mathbf{M}/\mathbf{D})^{-1}\mathbf{B}\mathbf{D}^{-1}\end{bmatrix}
	\end{equation}
\end{remark}

\begin{remark}[Combining both formulations]
	When both $\mathbf{A}$ and $\mathbf{D}$ are invertible, the Woodbury lemma proves that the two formulations give the same matrices for each block. In particular, they can be combined block by block at will.
\end{remark}

\section{Woodbury Inversion Formula}

\begin{proposition}[Inversion identity]\label{th:simplified_Woodbury}
	For $n\in\mathbb{N}$, let $\mathbf{P}\in\mathcal{M}_{n,n}(\mathbb{R})$ such that $\mathbf{I}+\mathbf{P}$ is invertible. Then:
	\begin{equation}
		(\mathbf{I}+\mathbf{P})^{-1}=\mathbf{I}-(\mathbf{I}+\mathbf{P})^{-1}\mathbf{P}
	\end{equation}
\end{proposition}

\begin{proof}
	This result is direct from:
	\begin{equation}
		\mathbf{I}=(\mathbf{I}+\mathbf{P})^{-1}(\mathbf{I}+\mathbf{P})=(\mathbf{I}+\mathbf{P})^{-1}+(\mathbf{I}+\mathbf{P})^{-1}\mathbf{P}
	\end{equation}
\end{proof}

\begin{proposition}[Push-through identity \cite{bib:matrix-cookbook}]\label{th:push_through}
	For $n,p\in\mathbb{N}$, let $\mathbf{P}\in\mathcal{M}_{n,p}(\mathbb{R})$ and $\mathbf{Q}\in\mathcal{M}_{p,n}(\mathbb{R})$ such that $\mathbf{I}_n+\mathbf{P}\mathbf{Q}$ is invertible. Then $\mathbf{I}_p+\mathbf{Q}\mathbf{P}$ is also invertible and:
	\begin{equation}
		(\mathbf{I}_n+\mathbf{P}\mathbf{Q})^{-1}\mathbf{P}=\mathbf{P}(\mathbf{I}_p+\mathbf{Q}\mathbf{P})^{-1}
	\end{equation}
\end{proposition}

\begin{proof}
	We start by using the identity:
	\begin{equation}
		\mathbf{P}+\mathbf{P}\mathbf{Q}\mathbf{P}=\mathbf{P}(\mathbf{I}_p+\mathbf{Q}\mathbf{P})=(\mathbf{I}_n+\mathbf{P}\mathbf{Q})\mathbf{P}
	\end{equation}

	Thus:
	\begin{equation}
		\begin{aligned}
		&\mathbf{P}(\mathbf{I}_p+\mathbf{Q}\mathbf{P})=(\mathbf{I}_n+\mathbf{P}\mathbf{Q})\mathbf{P}\\
		\iff&\mathbf{P}=(\mathbf{I}_n+\mathbf{P}\mathbf{Q})\mathbf{P}(\mathbf{I}_p+\mathbf{Q}\mathbf{P})^{-1}\\
		\iff&(\mathbf{I}_n+\mathbf{P}\mathbf{Q})^{-1}\mathbf{P}=\mathbf{P}(\mathbf{I}_p+\mathbf{Q}\mathbf{P})^{-1}
		\end{aligned}
	\end{equation}
	with $\mathbf{I}_p+\mathbf{Q}\mathbf{P}$ admitting a right inverse, hence being invertible.
\end{proof}

\begin{proposition}[Generalised Woodbury formula \cite{bib:Woodbury}]\label{th:Woodbury}
	For $n,p\in\mathbb{N}$, let $\mathbf{A}\in\mathcal{M}_{n,n}(\mathbb{R})$, $\mathbf{B}\in\mathcal{M}_{n,p}(\mathbb{R})$, $\mathbf{C}\in\mathcal{M}_{p,n}(\mathbb{R})$, and $\mathbf{D}\in\mathcal{M}_{p,p}(\mathbb{R})$. We assume that $\mathbf{A}$, $\mathbf{D}$, and $\mathbf{D}+\mathbf{C}\mathbf{A}^{-1}\mathbf{B}$ are invertible. Then $\mathbf{A}+\mathbf{B}\mathbf{D}^{-1}\mathbf{C}$ is also invertible, and:
	\begin{equation}
		\left(\mathbf{A}+\mathbf{B}\mathbf{D}^{-1}\mathbf{C}\right)^{-1}=\mathbf{A}^{-1}-\mathbf{A}^{-1}\mathbf{B}\left(\mathbf{D}+\mathbf{C}\mathbf{A}^{-1}\mathbf{B}\right)^{-1}\mathbf{C}\mathbf{A}^{-1}
	\end{equation}
\end{proposition}

\begin{proof}
	We first factorise:
	\begin{equation}
		\begin{aligned}
		\left(\mathbf{A}+\mathbf{B}\mathbf{D}^{-1}\mathbf{C}\right)^{-1}&=\left(\mathbf{A}(\mathbf{I}+\mathbf{A}^{-1}\mathbf{B}\mathbf{D}^{-1}\mathbf{C})\right)^{-1}\\
		&=\left(\mathbf{I}+\mathbf{A}^{-1}\mathbf{B}\mathbf{D}^{-1}\mathbf{C}\right)^{-1}\mathbf{A}^{-1}
		\end{aligned}
	\end{equation}

	Applying \Cref{th:simplified_Woodbury} to $\mathbf{A}^{-1}\mathbf{B}\mathbf{D}^{-1}\mathbf{C}$ yields:
	\begin{equation}
		\begin{aligned}
		\left(\mathbf{A}+\mathbf{B}\mathbf{D}^{-1}\mathbf{C}\right)^{-1}&=\left(\mathbf{I}-\left(\mathbf{I}+\mathbf{A}^{-1}\mathbf{B}\mathbf{D}^{-1}\mathbf{C}\right)^{-1}\mathbf{A}^{-1}\mathbf{B}\mathbf{D}^{-1}\mathbf{C}\right)\mathbf{A}^{-1}\\
		&=\mathbf{A}^{-1}-\left(\mathbf{I}+\mathbf{A}^{-1}\mathbf{B}\mathbf{D}^{-1}\mathbf{C}\right)^{-1}\mathbf{A}^{-1}\mathbf{B}\mathbf{D}^{-1}\mathbf{C}\mathbf{A}^{-1}
		\end{aligned}
	\end{equation}

	Using \Cref{th:push_through} with $\mathbf{P}=\mathbf{A}^{-1}\mathbf{B}$ and $\mathbf{Q}=\mathbf{D}^{-1}\mathbf{C}$ further gives:
	\begin{equation}
		\left(\mathbf{I}+\mathbf{A}^{-1}\mathbf{B}\mathbf{D}^{-1}\mathbf{C}\right)^{-1}\mathbf{A}^{-1}\mathbf{B}=\mathbf{A}^{-1}\mathbf{B}\left(\mathbf{I}+\mathbf{D}^{-1}\mathbf{C}\mathbf{A}^{-1}\mathbf{B}\right)^{-1}
	\end{equation}

	Therefore:
	\begin{equation}
		\begin{aligned}
		\left(\mathbf{A}+\mathbf{B}\mathbf{D}^{-1}\mathbf{C}\right)^{-1}&=\mathbf{A}^{-1}-\mathbf{A}^{-1}\mathbf{B}\left(\mathbf{I}+\mathbf{D}^{-1}\mathbf{C}\mathbf{A}^{-1}\mathbf{B}\right)^{-1}\mathbf{D}^{-1}\mathbf{C}\mathbf{A}^{-1}\\
		&=\mathbf{A}^{-1}-\mathbf{A}^{-1}\mathbf{B}\left(\mathbf{D}(\mathbf{I}+\mathbf{D}^{-1}\mathbf{C}\mathbf{A}^{-1}\mathbf{B})\right)^{-1}\mathbf{C}\mathbf{A}^{-1}\\
		&=\mathbf{A}^{-1}-\mathbf{A}^{-1}\mathbf{B}\left(\mathbf{D}+\mathbf{C}\mathbf{A}^{-1}\mathbf{B}\right)^{-1}\mathbf{C}\mathbf{A}^{-1}
		\end{aligned}
	\end{equation}
\end{proof}
